及
\[
\mathbf{T}_2=\sum \mathbf{x}_j\mathbf{x}_j'=(n-1)\mathbf{S}+n\bar{\mathbf{x}}\,\bar{\mathbf{x}}'
\]
为依赖的. 在此种情形下, 算法的具体步骤如下(假定总体的均值 $\boldsymbol{\mu}$ 与方差 $\boldsymbol{\Sigma}$ 未知):

\mvbold{预测步骤}\quad 对每一具有缺损值的向量 $\mathbf{x}_j$, 记 $\mathbf{x}_j^{(1)}$ 为其缺损的分量, $\mathbf{x}_j^{(2)}$ 为其可获取得分量. 于是 $\mathbf{x}_j'=[\mathbf{x}_j^{(1)'},\mathbf{x}_j^{(2)'}]$.

设由估计步骤算出的估计值 $\hat{\boldsymbol{\mu}}$ 和 $\hat{\boldsymbol{\Sigma}}$ 给定, 在 $\mathbf{x}^{(2)}$ 给定的条件下, 利用 $\mathbf{x}^{(1)}$ 的条件正态分布的均值来估计缺损值, 即用\footnote{若所有分量 $\mathbf{x}_j$ 均缺损, 则令 $\hat{\mathbf{x}}_j=\hat{\boldsymbol{\mu}}$ 且 $\hat{\mathbf{x}}_j\hat{\mathbf{x}}_j'=\hat{\boldsymbol{\Sigma}}+\hat{\boldsymbol{\mu}}\hat{\boldsymbol{\mu}}'$.}
\phantomsection\label{eq:5-38}
\begin{equation}
\hat{\mathbf{x}}_j^{(1)}=E(\mathbf{X}_j^{(1)}\mid\mathbf{x}_j^{(2)};\hat{\boldsymbol{\mu}},\hat{\boldsymbol{\Sigma}})
=\hat{\boldsymbol{\mu}}^{(1)}+\hat{\boldsymbol{\Sigma}}_{12}\hat{\boldsymbol{\Sigma}}_{22}^{-1}(\mathbf{x}_j^{(2)}-\hat{\boldsymbol{\mu}}^{(2)}) \tag{5-38}
\end{equation}
来估计 $\mathbf{x}_j^{(1)}$ 对 $\mathbf{T}_1$ 的贡献.

然后预测出 $\mathbf{x}_j^{(1)}$ 对 $\mathbf{T}_2$ 的贡献为
\phantomsection\label{eq:5-39}
\begin{equation}
\hat{\mathbf{x}}_j^{(1)}\hat{\mathbf{x}}_j^{(1)'}=E(\mathbf{X}_j^{(1)}\mathbf{X}_j^{(1)'}\mid\mathbf{x}_j^{(2)};\hat{\boldsymbol{\mu}},\hat{\boldsymbol{\Sigma}})
=\hat{\boldsymbol{\Sigma}}_{11}-\hat{\boldsymbol{\Sigma}}_{12}\hat{\boldsymbol{\Sigma}}_{22}^{-1}\hat{\boldsymbol{\Sigma}}_{21}
+\hat{\mathbf{x}}_j^{(1)}\hat{\mathbf{x}}_j^{(1)'} \tag{5-39}
\end{equation}
和
\[
\widehat{\mathbf{x}_j^{(1)}\mathbf{x}_j^{(2)'}}=E(\mathbf{X}_j^{(1)}\mathbf{X}_j^{(2)'}\mid\mathbf{x}_j^{(2)};\hat{\boldsymbol{\mu}},\hat{\boldsymbol{\Sigma}})
=\hat{\mathbf{x}}_j^{(1)}\mathbf{x}_j^{(2)'}
\]
将\hyperref[eq:5-38]{式(5-38)}和\hyperref[eq:5-39]{式(5-39)}中的贡献对所有含缺损分量的 $\mathbf{x}_j$ 求和, 此和与样本数据结合便可产生 $\bar{\mathbf{T}}_1$ 和 $\bar{\mathbf{T}}_2$.

\mvbold{估计步骤}\quad 计算修正后的极大似然估计值(见结论 4.11)
\phantomsection\label{eq:5-40}
\begin{equation}
\hat{\boldsymbol{\mu}}=\frac{\bar{\mathbf{T}}_1}{n}, \quad
\hat{\boldsymbol{\Sigma}}=\frac{1}{n}\bar{\mathbf{T}}_2-\hat{\boldsymbol{\mu}}\hat{\boldsymbol{\mu}}' \tag{5-40}
\end{equation}
我们通过例 5.13 来说明预测-估计算法的计算步骤.

\noindent\textbf{例 5.13}\quad (EM 算法的说明)

\mvbold{用不完全数据集}
\[
\mathbf{X}=\begin{bmatrix}
-&0&3\\
7&2&6\\
5&1&2\\
-&-&5
\end{bmatrix}
\]
来估计正态总体均值 $\boldsymbol{\mu}$ 与协方差阵 $\boldsymbol{\Sigma}$, 这里 $n=4$, $p=3$. 观测值向量 $\mathbf{x}_1$ 与 $\mathbf{x}_4$ 的某些分量的数据缺损.

我们先从可获观测值计算出样本均值初始估计:
\[
\hat{\mu}_1=\frac{7+5}{2}=6, \quad
\hat{\mu}_2=\frac{0+2+1}{3}=1, \quad
\hat{\mu}_3=\frac{3+6+2+5}{4}=4
\]
用这些均值代替缺损值, 比如说 $\hat{x}_{11}=6$, 我们就能得到协方差的初始估计, 在构造这些估计时, 均需除以 $n$, 因为这种算法最终产生的是极大似然估计 $\hat{\boldsymbol{\Sigma}}$. 于是
\[
\hat{\sigma}_{11}=\frac{(6-6)^2+(7-6)^2+(5-6)^2+(6-6)^2}{4}=\frac{1}{2}
\]
\[
\hat{\sigma}_{22}=\frac{1}{2}, \quad \hat{\sigma}_{33}=\frac{5}{2}
\]
